% Sección: Análisis exploratorio de datos

Los datos completos cuentan con 569 observaciones y 31 variables, de las cuales
una, \verb|target|, es la variable de respuesta que nos indica si el tumor al
que corresponde cada observación es maligno (\verb|M|) o benigno (\verb|B|), y
las demás son variables numéricas predictoras que resumen la información
obtenida de las células de dichos tumores. Como se mencionó en la introducción,
estas variables corresponden a tres medidas resumen (la media, el desvío estándar
y el peor caso) de diez características celulares listadas a continuación:

\begin{itemize}
  \item \verb|radius|: El radio de las células
  \item \verb|texture|: Variación en la intensidad del núcleo
  \item \verb|perimeter|: Perímetro del núcleo
  \item \verb|area|: Área del núcleo
  \item \verb|smoothness|: Variación local del radio de las células
  \item \verb|compactness|: Qué tan compacta o extendida es la célula
  \item \verb|concavity|: Profundidad de las partes cóncavas del contorno de las células
  \item \verb|concave|: Número de secciones cóncavas del contorno de las células
  \item \verb|symmetry|: Qué tan simétrica es la forma de las células
  \item \verb|fractal_dimension|: Complejidad del borde de las células
\end{itemize}

Antes de cualquier partición de los datos, observamos que la proporción de
observaciones de tumores malignos es 0.3726 aproximadamente (212
observaciones); es decir, los datos están ligeramente desbalanceados. Sin
embargo, dado que el desbalance no es tan marcado, se optó por no usar estratos
en la partición en datos de entrenamiento y testeo para minimizar las
decisiones tomadas sobre la totalidad de los datos. Se retuvo aproximadamente
el 30\% de las observaciones (171) para la evaluación de los modelos y el resto
se utilizó para el análisis exploratorio, el entrenamiento y la validación.

En los datos de entrenamiento la proporción de tumores malignos es 0.3643
aproximadamente. Afortunadamente no hay datos perdidos en ninguna de las
variables. En el Cuadro \ref{tab:resumen} se muestran algunas medidas resumen
de las variables predictoras, agrupadas según la medida resumen (valga la
redundancia) a la que correspondan. Las variables predictoras tienen unidades sumamente
dispares; por ejemplo, \verb|concave_mean| tiene valores entre 0 y 0.19
secciones cóncavas mientras que \verb|area_mean| los tiene entre 143 y 2501
unidades de área. Esta información es relevante ya que querríamos evitar que
algunas características tengan mayor impacto en el entrenamiento de los modelos
que otras. Todas las variables tienen rangos positivos y las posiciones de sus
medias y medianas sugieren que la mayoría son asimétricas con un sesgo hacia
los valores más altos.

Las Figuras \ref{fig:area}-\ref{fig:smooth} y
\ref{fig:radius}-\ref{fig:fractal}\footnote{La mayoría de los gráficos
exploratorios se presentan en el Anexo \ref{sec:anexo2} ya que son numerosos y
ocupan bastante espacio} ofrecen visualizaciones de las distribuciones de las
variables del conjunto de datos. Cada figura corresponde a una característica
de las células; la fila superior muestra histogramas de densidad según el valor
de \verb|target| para cada medida resumen y la fila inferior muestra un
diagrama de dispersión por cada combinación de dos de esas medidas con los
puntos coloreados según los valores de \verb|target|.

% Gráficos area
\begin{figure}[h]
\includegraphics[width=\textwidth, center]{figuras/distribuciones\_area.pdf}
\caption{Histogramas de densidad de las variables asociadas a \textit{area} según \textit{target} (fila superior) y diagramas de dispersión de las mismas variables de a pares (fila inferior)}
\label{fig:area}
\end{figure}

Si nos enfocamos en los histogramas, se puede observar que, efectivamente,
algunas variables tienen marcados sesgos hacia los valores más altos, en
particular las que corresponden a los desvíos estándar de las características
de los tumores. Estas variables también tienden a presentar valores extremos;
esto se puede apreciar en la Figura \ref{fig:area}, por ejemplo, en el desvío
estándar del área de las células. También se observa que en casi todas las
variables los tumores malignos tienden a tener valores más altos. Las
excepciones son las variables relacionadas con la dimensión fractal del
contorno de las células y el desvío estándar de la simetría, la suavidad y la
textura. En algunos casos, como el área promedio de las células, las
distribuciones condicionales a la calidad de benigno o maligno del tumor son
más simétricas que la distribución general de la variable.

Un aspecto relevante de las variables es su dependencia; es de esperar que las
que corresponden a las mismas características de las células estén relacionadas
entre sí. Si examinamos los diagramas de dispersión de las Figuras
\ref{fig:area}-\ref{fig:smooth} y \ref{fig:radius}-\ref{fig:fractal} vemos que,
efectivamente, tienden a estar fuertemente asociadas con un patrón bastante
lineal. La correlación es más fuerte entre la media y el peor caso, con el
coeficiente $\rho$ de pearson yendo de 0.688 para la simetría hasta 0.9718 para
el perimetro. Es más débil entre la media y el desvío, yendo de 0.3505 para la
suavidad hasta 0.8077 para el área, y entre el desvío y el peor caso, de 0.4035
para la suavidad hasta 0.86 para el área. 

\begin{figure}[h]
\includegraphics[width=\textwidth, center]{figuras/distribuciones\_smoothness.pdf}
\caption{Histogramas de densidad de las variables asociadas a \textit{smoothness} según \textit{target} (fila superior) y diagramas de dispersión de las mismas variables de a pares (fila inferior)}
\label{fig:smooth}
\end{figure}

Si nos enfocamos en la condición de maligno o benigno de los tumores en los
diagramas de dispersión, vemos que, a diferencia de lo que se ve en los
histogramas, hay fronteras bastante nítidas entre ambos tipos en algunas de las
características. Esto es muy claro en las variables asociadas al radio (Figura
\ref{fig:radius}), el perímetro (Figura \ref{fig:perimeter}) y el área (Figura
\ref{fig:area}). En otros casos, como en las variables de suavidad (Figura
\ref{fig:smooth}), no parece haber una frontera muy clara. Nada de esto
descarta que sí exista una frontera nítida en dimensiones más altas; las
visualizaciones que se hicieron no van más allá de las dos dimensiones y,
teniendo 30 variables predictoras, es probable que no estemos notando patrones
importantes.
